\documentclass[10pt,a4paper]{amsart}
\usepackage{amsmath,amssymb,amsthm,graphicx}
\usepackage{verbatim}
\newenvironment{reference}{\comment}{\endcomment}
%\newenvironment{reference}{}{}
\newenvironment{slogan}{\comment}{\endcomment}
\newenvironment{history}{\comment}{\endcomment}
% For commutative diagrams we use Xy-pic
\usepackage[all]{xy}

% We use 2cell for 2-commutative diagrams.
\xyoption{2cell}
% Theorem environments.
%
\theoremstyle{plain}
\newtheorem{theorem}[subsection]{Theorem}
\newtheorem{proposition}[subsection]{Proposition}
\newtheorem{lemma}[subsection]{Lemma}

\theoremstyle{definition}
\newtheorem{definition}[subsection]{Definition}
\newtheorem{example}[subsection]{Example}
\newtheorem{exercise}[subsection]{Exercise}
\newtheorem{situation}[subsection]{Situation}

\theoremstyle{remark}
\newtheorem{remark}[subsection]{Remark}
\newtheorem{remarks}[subsection]{Remarks}

\numberwithin{equation}{subsection}

% Macros
%
\def\lim{\mathop{\mathrm{lim}}\nolimits}
\def\colim{\mathop{\mathrm{colim}}\nolimits}
\def\Spec{\mathop{\mathrm{Spec}}}
\def\Hom{\mathop{\mathrm{Hom}}\nolimits}
\def\Ext{\mathop{\mathrm{Ext}}\nolimits}
\def\SheafHom{\mathop{\mathcal{H}\!\mathit{om}}\nolimits}
\def\SheafExt{\mathop{\mathcal{E}\!\mathit{xt}}\nolimits}
\def\Sch{\mathit{Sch}}
\def\Mor{\mathop{\mathrm{Mor}}\nolimits}
\def\Ob{\mathop{\mathrm{Ob}}\nolimits}
\def\Sh{\mathop{\mathit{Sh}}\nolimits}
\def\NL{\mathop{N\!L}\nolimits}
\def\CH{\mathop{\mathrm{CH}}\nolimits}
\def\proetale{{pro\text{-}\acute{e}tale}}
\def\etale{{\acute{e}tale}}
\def\QCoh{\mathit{QCoh}}
\def\Ker{\mathop{\mathrm{Ker}}}
\def\Im{\mathop{\mathrm{Im}}}
\def\Coker{\mathop{\mathrm{Coker}}}
\def\Coim{\mathop{\mathrm{Coim}}}

% Boxtimes
%
\DeclareMathSymbol{\boxtimes}{\mathbin}{AMSa}{"02}

%
% Macros for moduli stacks/spaces
%
\def\QCohstack{\mathcal{QC}\!\mathit{oh}}
\def\Cohstack{\mathcal{C}\!\mathit{oh}}
\def\Spacesstack{\mathcal{S}\!\mathit{paces}}
\def\Quotfunctor{\mathrm{Quot}}
\def\Hilbfunctor{\mathrm{Hilb}}
\def\Curvesstack{\mathcal{C}\!\mathit{urves}}
\def\Polarizedstack{\mathcal{P}\!\mathit{olarized}}
\def\Complexesstack{\mathcal{C}\!\mathit{omplexes}}
% \Pic is the operator that assigns to X its picard group, usage \Pic(X)
% \Picardstack_{X/B} denotes the Picard stack of X over B
% \Picardfunctor_{X/B} denotes the Picard functor of X over B
\def\Pic{\mathop{\mathrm{Pic}}\nolimits}
\def\Picardstack{\mathcal{P}\!\mathit{ic}}
\def\Picardfunctor{\mathrm{Pic}}
\def\Deformationcategory{\mathcal{D}\!\mathit{ef}}
\usepackage{cleveref}
\crefname{originallemma}{original lemma}{original lemmas}
\Crefname{originallemma}{Original lemma}{Original lemmas}
\makeatletter
\newlabel{lemma-closed-immersion-proper}{{lemma-closed-immersion-proper}{source}{}{sourcepointer-1}{}}
\newlabel{lemma-closed-immersion-proper@cref}{{[originallemma][0][]lemma-closed-immersion-proper}{[1][0][]source}{}{}{}}
\makeatother
\crefname{originallemma}{original lemma}{original lemmas}
\Crefname{originallemma}{Original lemma}{Original lemmas}
\makeatletter
\newlabel{lemma-every-point-residual-gerbe-locally-Noetherian}{{lemma-every-point-residual-gerbe-locally-Noetherian}{source}{}{sourcepointer-2}{}}
\newlabel{lemma-every-point-residual-gerbe-locally-Noetherian@cref}{{[originallemma][0][]lemma-every-point-residual-gerbe-locally-Noetherian}{[1][0][]source}{}{}{}}
\makeatother
\crefname{originallemma}{original lemma}{original lemmas}
\Crefname{originallemma}{Original lemma}{Original lemmas}
\makeatletter
\newlabel{lemma-immersion-locally-finite-type}{{lemma-immersion-locally-finite-type}{source}{}{sourcepointer-3}{}}
\newlabel{lemma-immersion-locally-finite-type@cref}{{[originallemma][0][]lemma-immersion-locally-finite-type}{[1][0][]source}{}{}{}}
\makeatother
\crefname{originallemma}{original lemma}{original lemmas}
\Crefname{originallemma}{Original lemma}{Original lemmas}
\makeatletter
\newlabel{lemma-locally-finite-type-locally-noetherian}{{lemma-locally-finite-type-locally-noetherian}{source}{}{sourcepointer-4}{}}
\newlabel{lemma-locally-finite-type-locally-noetherian@cref}{{[originallemma][0][]lemma-locally-finite-type-locally-noetherian}{[1][0][]source}{}{}{}}
\makeatother
\crefname{originalreference}{original reference}{original references}
\Crefname{originalreference}{Original reference}{Original references}
\makeatletter
\newlabel{stacks-properties-lemma-reduced-closed-substack}{{stacks-properties-lemma-reduced-closed-substack}{source}{}{sourcepointer-5}{}}
\newlabel{stacks-properties-lemma-reduced-closed-substack@cref}{{[originalreference][0][]stacks-properties-lemma-reduced-closed-substack}{[1][0][]source}{}{}{}}
\makeatother
\title{The residual space of a point is closed in the space of the stack}
\author{The Stacks Project Authors}
\date{stacks-project:stacks-morphisms.tex:local-snapshot-20261012; GNU FDL 1.2 or later}
\begin{document}
\maketitle
\noindent\textit{Editorial note.} All delivered mathematical text is the author's own, copied from the pinned Stacks Project chapter source sources/stacks-algebraic-stacks-morphisms/pinned-original-source.tex; no statement or proof step was rewritten, reordered or removed.\par\medskip
\noindent\textit{Editorial note.} Preamble provenance: the Stacks Project keeps one single shared preamble for every chapter; this card takes it directly from the staged sources/stacks-algebraic-stacks-morphisms/preamble.tex as three verbatim spans (lines 8-12, 14-18 and 154-217), so no derived file was written and nothing was added or redefined. The wrapper class and the base packages are supplied by the assembler.\par\medskip
\noindent\textit{Editorial note.} Bibliography provenance: neither the counted statement nor its proof contains a citation (this lane rescans the whole span before writing the card), so the card carries no bibliography and no entry was reconstructed.\par\medskip
\noindent\textit{Editorial note.} Selection provenance: this card was opened under the file-based control contract (CONTROL.md), with the mechanical pre-card gate run on its statement and proof interval before the card was written.\par\medskip
\noindent\textit{Allowed standard prerequisites (editorial).}\begin{itemize}
\item The description of the residual gerbe of a point through the closure of the point. Used at the start of the counted proof (source lines 10857-10862). Identifies the underlying space of the residual gerbe with the closure of the point.
\item The statement that the closure of a point of the space of a stack is a closed subset whose preimage in a presentation is the closure of the corresponding point. Invoked inside the counted proof (source lines 10863-10868). Turns closedness of the residual space into a statement about a presentation.
\item The corresponding statement for schemes that the closure of a point is closed. Used at the end of the counted proof (source lines 10869-10870). Gives closedness after passing to a presentation.
\end{itemize}\medskip
\ifdefined\hypertarget\hypertarget{sourcepointer-1}{}\fi
\noindent\textit{Source reference.} Original lemma lemma-closed-immersion-proper: sources/stacks-algebraic-stacks-morphisms/pinned-original-source.tex, source lines 8635--8635. Navigation pointer only: the counted text refers to the Stacks item lemma-closed-immersion-proper, which is an existing labeled result of this same chapter and is not part of this card. The pointer carries the target's own label as its number and its definition line as its span; its content is not reproduced here.\par\medskip
\ifdefined\hypertarget\hypertarget{sourcepointer-2}{}\fi
\noindent\textit{Source reference.} Original lemma lemma-every-point-residual-gerbe-locally-Noetherian: sources/stacks-algebraic-stacks-morphisms/pinned-original-source.tex, source lines 7165--7165. Navigation pointer only: the counted text refers to the Stacks item lemma-every-point-residual-gerbe-locally-Noetherian, which is an existing labeled result of this same chapter and is not part of this card. The pointer carries the target's own label as its number and its definition line as its span; its content is not reproduced here.\par\medskip
\ifdefined\hypertarget\hypertarget{sourcepointer-3}{}\fi
\noindent\textit{Source reference.} Original lemma lemma-immersion-locally-finite-type: sources/stacks-algebraic-stacks-morphisms/pinned-original-source.tex, source lines 3514--3514. Navigation pointer only: the counted text refers to the Stacks item lemma-immersion-locally-finite-type, which is an existing labeled result of this same chapter and is not part of this card. The pointer carries the target's own label as its number and its definition line as its span; its content is not reproduced here.\par\medskip
\ifdefined\hypertarget\hypertarget{sourcepointer-4}{}\fi
\noindent\textit{Source reference.} Original lemma lemma-locally-finite-type-locally-noetherian: sources/stacks-algebraic-stacks-morphisms/pinned-original-source.tex, source lines 3525--3525. Navigation pointer only: the counted text refers to the Stacks item lemma-locally-finite-type-locally-noetherian, which is an existing labeled result of this same chapter and is not part of this card. The pointer carries the target's own label as its number and its definition line as its span; its content is not reproduced here.\par\medskip
\ifdefined\hypertarget\hypertarget{sourcepointer-5}{}\fi
\noindent\textit{Source reference.} Original reference stacks-properties-lemma-reduced-closed-substack: sources/stacks-algebraic-stacks-morphisms/pinned-original-source.tex, source lines 6854--6854. Navigation pointer only: the counted text refers to the Stacks item stacks-properties-lemma-reduced-closed-substack, which belongs to another chapter of the project and is not part of this card. The pointer carries the target's own label as its number because this chapter's staged snapshot does not contain the other chapter's tag table; its content is not reproduced here.\par\medskip
\begin{lemma}
\label{lemma-residual-space-closed}
Let $\mathcal{X}$ be a locally Noetherian algebraic stack.
Let $x \in |\mathcal{X}|$ with residual gerbe
$\mathcal{Z}_x \subset \mathcal{X}$
(Lemma \ref{lemma-every-point-residual-gerbe-locally-Noetherian}).
Then $x$ is a closed point of $|\mathcal{X}|$ if and only if
the morphism $\mathcal{Z}_x \to \mathcal{X}$ is a closed immersion.
\end{lemma}

\begin{proof}
If $\mathcal{Z}_x \to \mathcal{X}$ is a closed immersion, then
$x$ is a closed point of $|\mathcal{X}|$, see for example
Lemma \ref{lemma-closed-immersion-proper}.
Conversely, assume $x$ is a closed point of $|\mathcal{X}|$.
Let $\mathcal{Z} \subset \mathcal{X}$
be the reduced closed substack with $|Z| = \{x\}$
(Properties of Stacks,
Lemma \ref{stacks-properties-lemma-reduced-closed-substack}).
Then $\mathcal{Z}$ is a locally Noetherian algebraic stack
by Lemmas \ref{lemma-immersion-locally-finite-type} and
\ref{lemma-locally-finite-type-locally-noetherian}.
Since also $\mathcal{Z}$ is reduced and $|\mathcal{Z}| = \{x\}$
it follows that $\mathcal{Z} = \mathcal{Z}_x$ is the
residual gerbe by definition.
\end{proof}

\end{document}
